\section{Disputed questions and the Byzantine witness}

\subsection{Bonaventure: composition without a body}

Bonaventure and Thomas agree that an angel is a creature, intellectually alive,
free, and incorporeal. They disagree over whether its created substance is
composed of matter and form. Bonaventure's question concerns the principles
that make a finite spirit a recipient of perfection. Mutability, receptivity,
individuality, and limitation seem to him to require an actual and a potential
principle within its essence. He therefore judges composition from matter and
form the truer account.\footnote{Bonaventure, \emph{In II Sententiarum},
d.~3, p.~1, a.~1, q.~1, response; \emph{Opera omnia}, vol.~2
(Quaracchi, 1885), pp.~90--91;
\url{https://archive.org/details/doctorisseraphic02bona/page/90/mode/2up}.}

His use of \emph{matter} cannot simply be replaced with \emph{body}. At the
opening of the next question he explicitly takes the term broadly for a
potential principle entering into a thing's constitution. This allows him
to speak of spiritual matter without giving angels weight, bodily organs,
or quantitative extension. In the \emph{Breviloquium}, he accordingly calls
the angelic substance simple, spiritual, and incorporeal, lacking quantitative
dimensions. These descriptions do not retract his account of composition:
they concern a different kind of simplicity.\footnote{Bonaventure,
\emph{In II Sententiarum}, d.~3, p.~1, a.~1, q.~2, opening statement,
\emph{Opera omnia}, vol.~2, p.~94;
\emph{Breviloquium}, II.6, especially the explanation of the angelic
substance's simplicity; Latin text,
\url{https://catholiclibrary.org/library/view?docId=/Medieval-OR/BonaventuraSBreviloquium.00000158.la.html\&chunk.id=00000039}.}

Thomas rejects matter in an intellectual substance because intellectual
reception differs from material reception. Matter receives a form by becoming
the kind of thing that form determines; an intellect can know many forms
without successively becoming the physical things it knows. But Thomas does
not therefore make an angel pure actuality. Its essence receives existence;
its finite being depends upon God. Both theologians defend the creature's
dependence. Their disagreement concerns whether that dependence and capacity
require matter within the spiritual substance.\footnote{\ST{I, q.~50,
a.~2, corpus and ad~2--4};
\url{https://www.newadvent.org/summa/1050.htm}.}

Their two accounts articulate the same fundamental distance from divine
simplicity through different metaphysical principles. For Bonaventure,
created receptivity reaches into the angel's substantial constitution as
matter and form. For Thomas, the subsistent form itself receives existence.
The angel's immateriality therefore magnifies the goodness of its Creator
without erasing the difference between having being and being its source.

\subsection{John of Damascus: created light and finite presence}

John of Damascus describes angels as intelligent beings created through the
Word and perfected through the Spirit. Their life is received; their
illumination comes from God. His declaration that only God is absolutely
immaterial places even the highest spirit at an immeasurable distance from
the Creator. He nevertheless denies bodily dimensions to angels. Their
incorporeality does not give them divine omnipresence: their activity is
limited to the place of their commission. Their visible manifestations
accommodate human perception rather than exposing their essence.\footnote{John
of Damascus, \emph{Exposition of the Orthodox Faith}, II.3, trans.
S.~D.~F. Salmond, \emph{Nicene and Post-Nicene Fathers}, second
series, vol.~9, New Advent witness,
\url{https://www.newadvent.org/fathers/33042.htm}.}

This language should not be turned into an identity with either later
scholastic account. John does not formulate Bonaventure's question about
spiritual matter. He also leaves undecided whether angels share one essence
or differ in essence. Thomas, by contrast, determines that distinct
immaterial angels differ specifically. John's reserve and Thomas's argument
are different answers to how far the theologian can describe created
intellectual diversity.\footnote{John of Damascus, \emph{Exposition}, II.3,
paragraph beginning ``Whether they are equals in essence''; \ST{I, q.~50,
a.~4}.}

Their disagreement about creation is still more explicit. John prefers the
creation of the intellectual world before the visible world, following
Gregory the Theologian. Thomas regards simultaneous creation as more probable:
angels and bodies are parts of one ordered universe. He expressly declines
to condemn the contrary view, acknowledging the authority of Gregory.
Angelic creation from nothing is the doctrinal ground; its priority relative
to bodily creation remains \textbf{disputed}.\footnote{John of Damascus,
\emph{Exposition}, II.3, final discussion of the time of creation;
\ST{I, q.~61, a.~3, corpus and replies},
\url{https://www.newadvent.org/summa/1061.htm}.}

The same Byzantine witness makes the fall an act of freedom, not a defect
planted in nature by its Maker. The first rebel was created good and departed
from that goodness. God's permission bounds the demons' activity; they do
not acquire the Creator's knowledge of the future. Here John's account
supports the distinction between a spirit's natural gifts and its moral
condition. Intelligence can remain while its exercise becomes hostile to
its proper end.\footnote{John of Damascus, \emph{Exposition}, II.4.}

\subsection{The first sin: pride, envy, and the object desired}

Augustine places the angels' happiness in their enjoyment of the unchangeable
good and their certain possession of it forever. The devil's rebellion turns
from that participation toward a private lordship. Such a lordship cannot
escape the dominion of the Almighty: the creature's proud claim contradicts
its actual dependence. The evil lies in the will that abandons truth, while
the angelic nature remains the work of God. Augustine's reading of the devil
as a sinner ``from the beginning'' consequently does not make evil his
created essence. It identifies the beginning of his sin with the pride
through which he departed from the good.\footnote{Augustine, \emph{City of
God}, XI.13--15, trans. Marcus Dods;
\url{https://www.newadvent.org/fathers/120111.htm}.}

Thomas argues that an intellectual
creature's first disorder must be pride: a spiritual good is desired without
the subjection owed to God's rule. Envy follows when another creature's good
appears to threaten the singular excellence the rebel seeks. Desire for
equality with God cannot mean a reasoned expectation of becoming the divine
essence. Thomas locates the disorder instead in seeking beatitude through
one's own power or according to one's own measure.\footnote{\ST{I, q.~63,
aa.~2--3}; \url{https://www.newadvent.org/summa/1063.htm}.}

Bonaventure's account gives this self-exaltation a closely related form.
Lucifer presumes upon his own excellence and desires precedence over others.
By resting in his private good he treats himself as his own first principle
and final good. The angel cannot actually become either. The attempted
ascent is therefore a fall, and subsequent hostility toward humanity
expresses the disordered will.\footnote{Bonaventure,
\emph{Breviloquium}, II.7, opening and explanatory paragraphs; Latin text,
\url{https://catholiclibrary.org/library/view?docId=/Medieval-OR/BonaventuraSBreviloquium.00000158.la.html\&chunk.id=00000041}.}

The identity and original rank of the chief rebel remain a separate question.
John places him over the earthly realm. Thomas allows that position without
judging it contrary to faith, but finds it more probable that the highest
angel fell: greater excellence could supply a greater occasion for pride,
even though it also gave less natural inclination to sin. A greater gift
can become the object of greater presumption. Neither account makes the gift
itself evil.\footnote{John of Damascus, \emph{Exposition}, II.4;
\ST{I, q.~63, a.~7}.}

In these accounts, the fall reveals the difference between possessing a gift
and receiving it rightly. The angel's power and intelligence are good. Pride
turns their possessor from thanksgiving to self-sufficiency; envy then
opposes the good shared with another. Angelic holiness makes the opposite
movement: the spirit receives its excellence from God and delights in the
communication of divine goodness.
